\section{Experimental Validation}
\label{sec:hardexp}
\subsection{Available Measurements}
The preliminary directional test is retained in Appendix~\ref{app:poc}. It compares eight angular positions at $0.30\,\mathrm m$ and supports a coarse orientation-dependent response. It does not establish a local Gaussian slope at the operating tone, range accuracy over the tracking domain, or the paired-measurement update rate.

The supplied static recording dated 22 October 2026 contains five packets with median pair period $12.84\,\mathrm s$, zero accepted measurements, and zero nonzero commands in offline replay. Three silence recordings contain two packets each and no valid measurements. These records demonstrate the current data limitation and the software's refusal to move on unusable data. They do not demonstrate closed-loop tracking.

\subsection{Hardware Identification and Sensing}
\pending{E1 --- Bench identification}{Insert final sensor calibration and actuator identification from H1--H4. Give trial dates, firmware/configuration versions, sample counts, training versus held-out split, reference instruments, and measurement uncertainties. Include distance/angle grids, source orientation, motor-on noise, and pair-timing distributions.}

\subsection{Moving-Target Following}
\pending{E2 --- Straight following and speed change}{Run repeatable marked trajectories with independently recorded target/follower references. Report range/bearing errors, settling/overshoot with stated time origin, control effort, saturation, TRACK/HOLD occupancy, minimum clearance, timing, and every failed trial. Compare LQI, pole placement, and the shared-supervisor baseline using the same calibrated limits.}

\pending{E3 --- Turning and Signal Loss}{Measure feasible turning trajectories, a controlled source interruption, reacquisition delay, and the effect of target/speaker aspect. Include synchronized trajectory plots and logs. Keep the target within demonstrated speed and curvature limits.}

\subsection{Stationary Approach and Model Agreement}
\pending{E4 --- Terminal approach}{Test the explicitly selected stationary mission separately from moving-target regulation. Report terminal separation distribution, braking distance, stop latch behavior, and motor-deadband effects.}

\pending{E5 --- Model versus hardware}{Compare measured small-perturbation trajectories with the local LTI prediction and the identified nonlinear simulation, using matched initial conditions and commands. Quantify prediction errors and the domain over which the model is useful. A tape-marked path needs a timestamped position/heading reference; an IMU alone is not position ground truth.}

\section{Discussion and Limitations}
The nominal design is controllable and observable only within its moving operating regime. At rest, loss of heading information coincides with loss of steering authority. The separate terminal mission avoids claiming that the same LQI equilibrium applies to both moving pursuit and stationary arrival.

The principal unresolved limitation is acoustic range identifiability. A multiplicative source or receiver gain change is indistinguishable from range under Eq.~\eqref{eq:amp}. The simulated bias test exposes this ambiguity directly. Recalibration and fixed source settings can reduce drift, but only held-out measurements can establish an acceptable operating domain. If amplitude range is not repeatable, the project needs an additional observable quantity or a revised task; controller retuning alone is insufficient.

Sequential sensing introduces another constraint. The two amplitudes correspond to different positions and orientations. Covariance inflation accounts for bounded mismatch, while delayed replay accounts for measurement age; neither reconstructs information lost through an invalid measurement model. The fixed-rate control loop also runs on a non-real-time operating system. Measured scheduling delay, clock drift, and transport uncertainty must therefore accompany any claim about bandwidth or recovery time.

The observer predicts follower motion without encoder feedback. Motor deadband, slip, battery voltage, and load can bias that prediction. Existing actuator stress tests cover only selected scale and lag errors. Similarly, synthetic reflection events do not cover persistent multipath, several sources, non-Gaussian ambiguity, or source occlusion. No global convergence or arbitrary-obstacle guarantee is asserted.

The demonstrated contribution at this stage is a reproducible connection between a relative state-space model, realizable front-steered feedback, delayed acoustic estimation, and explicit failure tests. The experimental contribution remains open until sensing accuracy, actuator behavior, and closed-loop following have been measured.

\pending{D1 --- Discussion from measured evidence}{Explain where the measured model agrees or fails, whether integral control improves distance regulation over the baseline, which constraint limits bandwidth, and which conclusions survive changes in source aspect and room conditions. Avoid generalizing beyond the tested domain.}

\section{Conclusion}
The moving-target local model admits observable and controllable state feedback, whereas the stopped vehicle requires a separate terminal treatment. Discrete LQI and channel-wise pole placement regulate simulated following distance under the nominal conditions and selected disturbances; a constant-speed bearing steering law does not. Calibration bias and slow sequential capture remain decisive limitations. The software and simulation evidence supports proceeding to gated hardware trials, but does not yet establish a successful physical tracking system.

\pending{C1 --- Final experimental conclusion}{Replace this provisional conclusion after E1--E5: state measured tracking/terminal accuracy, operating speed and sensing rate, comparative controller results, failed objectives, and the validated operating domain. Update the abstract consistently.}
